\domainevide{Les cotes de ce thème supposent les schémas abéliens et le foncteur de Picard, absents de mathlib.}{9, 14, 44, 47, 49, 50, 51, 52, 57, 59, 100, 130, 139}
\noindent Pistes de l'onglet Findings pour ces cotes, non démontrées ici :
\begin{itemize}\small\setlength{\itemsep}{0pt}
\item Cote~9 : pour un schéma abélien ordinaire localement polarisable sur une base connexe de caractéristique $p$, $\det T_p(\tilde\varphi)_{\text{ét}}/\sqrt{\deg\tilde\varphi}$ ne dépend pas de la polarisation et donne un isomorphisme $\det T_p(A)_{\text{ét}} \simeq \det T_p(A^*)_{\text{ét}}$. Mathlib n'a pas les schémas abéliens ni les modules de Tate $p$-adiques.\piste{9-det-etale-tate-iso}
\item Cote~9 : pour une variété abélienne ordinaire polarisée et $u \in \mathrm{End}(A) \otimes \mathbf Q$ symétrique et positif pour l'involution de Rosati, le déterminant de $u$ sur la partie étale du module de Tate $p$-adique vaut $+\sqrt{\deg u}$. Mathlib n'a pas les variétés abéliennes ni l'involution de Rosati.\piste{9-symmetric-positive-det-sqrt-degree}
\item Cote~14 : une équivalence d'Albanese sur les cycles de toute codimension, comprise entre les équivalences rationnelle et algébrique ; dans l'anneau quotient, le produit de deux classes algébriquement triviales est nul. Mathlib n'a pas les cycles algébriques des variétés projectives lisses.\piste{14-albanese-equivalence-ring}
\item Cote~14 : un faisceau inversible $L$ sur un schéma projectif est ample dès que $\chi(\mathcal O_Z \otimes L^{\otimes n}) \to +\infty$ pour tout sous-schéma fermé intègre $Z$ de dimension $\geq 1$ (énoncé attribué à Mumford et Nakai). Mathlib n'a pas la caractéristique d'Euler des faisceaux cohérents ni l'amplitude.\piste{14-nakai-euler-characteristic-criterion}
\item Cote~44 : une page biffée sur les systèmes de racines semble déduire l'intégralité des nombres de Cartan de trois axiomes (finitude, réduction, symétries) ; ces trois axiomes seuls ne l'entraînent pas. Accessible : un contre-exemple d'algèbre linéaire sur $\mathbf Q$, $\Delta = \{\pm 3e_1, \pm 3e_2, \pm(e_1 \pm e_2)\}$, reste à écrire.\piste{44-cartan-integrality-from-axioms}
\item Cote~47 : sur une base $S$, Grothendieck--Riemann--Roch donne $\mathrm{ch}(x * y) = (\mathrm{ch}\,x * \mathrm{ch}\,y)\cdot\mathrm{Todd}(t_A)$ sur un schéma abélien $A$, et la transformation $\Psi_D$ est multiplicative sans correction. Mathlib n'a pas Grothendieck--Riemann--Roch ni les schémas abéliens.\piste{47-chow-fourier-todd-over-base}
\item Cote~47 : le dossier propose l'intégralité des classes $\binom{\chi}{k}\delta'^k/\chi^m$ sur une variété abélienne, comme conséquence de $c(\Phi_L(\mathrm{cl}\,L)) = (1 - D'/\chi)^\chi$. Mathlib n'a pas les classes de Chern sur les variétés abéliennes.\piste{47-dual-polarisation-divisibility}
\item Cote~49 : pour $\varphi : X \to Y$ de type fini sur $S$ localement noethérien, $F$ quasi-cohérent sur $X$ plat sur $S$ et $G$ cohérent sur $Y$, si $G \to \varphi_*F$ est injectif sur la fibre en $s$ aux points au-dessus de $y$, $G$ est plat sur $S$ en $y$. Mathlib n'a pas la platitude des faisceaux cohérents sur les schémas.\piste{49-flatness-lemma-pushforward}
\item Cote~50 : un cadre axiomatique pour les endomorphismes des variétés abéliennes (dualité, cônes, forme d'intersection), dont le degré $\nu(\alpha)$, les coefficients $\sigma_i(\alpha)$ et la formule des traces sont des conséquences, sans points ni diviseurs. Accessible une fois les axiomes posés ; les arguments de positivité allongent la preuve.\piste{50-axiomatic-trace-formula}
\item Cote~51 : si $X$ est projectif et plat sur un anneau de valuation discrète, à fibres lisses hors d'un fermé de codimension $\geq 2$, $\mathrm{Alb}(X_K)$ et $(\mathrm{Pic}^0_{X_K})_{\mathrm{red}}$ ont bonne réduction. Mathlib n'a pas les schémas d'Albanese et de Picard.\piste{51-albanese-good-reduction-codim-two}
\item Cote~57 : pour $X \subset \mathbf P^r$ normal intègre de dimension $\geq 2$ et sa section hyperplane générique $Y_{\bar K}$, $H^1(X, G) \to H^1(Y_{\bar K}, G)$ est injectif pour $G = \mathbf G_a$ et pour $G$ fini commutatif. Mathlib n'a pas la cohomologie plate des schémas.\piste{57-generic-hyperplane-h1-injective}
\item Cote~57 : deux sous-schémas en groupes fermés \q{rigides} d'un schéma en groupes de type fini sur une base connexe localement noethérienne, qui coïncident sur une fibre, sont égaux. Mathlib n'a pas les schémas en groupes sur une base noethérienne.\piste{57-rigid-subgroups-one-fibre}
\item Cote~57 : pour $f : X \to Y$ propre entre espaces analytiques complexes, $(R^1f_*\mathcal O^*_X)_y \to \varprojlim \mathrm{Pic}(X_n)$ est un isomorphisme si et seulement si $(R^1f_*\mathcal O_X)_y$ est de longueur finie. Mathlib n'a pas les espaces analytiques complexes ni GAGA.\piste{57-analytic-pic-completion-criterion}
\item Cote~59 : pour un schéma abélien $A/S$ de dimension relative $2$ et une section $\delta$ de $\mathrm{NS}_{A/S}$ avec $\chi(\delta)$ pair, le torseur $\mathrm{Pic}^\delta_{A/S}$ est trivial. Mathlib n'a pas les torseurs de Picard des schémas abéliens.\piste{59-surface-even-chi-torsor}
\item Cote~59 : le déterminant de la cohomologie du faisceau de Poincaré donne une trivialisation canonique d'un torseur formé d'$\mathrm{Alb}^{\chi(\delta)}_{X/S}$ et de deux torseurs de Picard. Mathlib n'a pas le déterminant de la cohomologie ni le faisceau de Poincaré.\piste{59-determinant-canonical-section}
\item Cote~100 : sur $S$ localement noethérien, $A \mapsto (A_U, A[p^\infty], \mathrm{id})$ est pleinement fidèle dès que l'ouvert $U$ rencontre chaque composante connexe de $S$, sans hypothèse de normalité ni de réduction. Mathlib n'a pas les schémas abéliens ni les groupes $p$-divisibles.\piste{100-full-faithfulness-connected-components}
\item Cote~100 : si $f : S' \to S$ est propre et $\mathcal O_S \to f_*\mathcal O_{S'}$ injectif, sans platitude, le changement de base le long de $f$ est pleinement fidèle sur les schémas abéliens. Mathlib n'a pas la descente des morphismes le long d'un morphisme propre.\piste{100-proper-epimorphism-hom-descent}
\item Cote~100 : un schéma abélien formel sur $\operatorname{Spf} R$ est algébrisable si et seulement si ses changements de base aux normalisés des quotients intègres de dimension un de $R$ le sont. Mathlib n'a pas l'algébrisation des schémas formels.\piste{100-formal-abelian-curve-criterion}
\item Cote~100 : dans le spectre d'un anneau local noethérien complet, une famille dénombrable de fermés dont la réunion contient tous les points de dimension $\leq 1$ a une sous-famille finie qui recouvre le spectre. Mathlib n'a pas l'indénombrabilité des idéaux premiers d'un anneau local complet.\piste{100-baire-dimension-one-points}
\item Cote~100 : sur un trait, tout homomorphisme de la fibre générique d'un groupe $p$-divisible vers celle d'un modèle de Néron se prolonge de façon unique, le théorème de Tate étant admis. Mathlib n'a pas les modèles de Néron ni les groupes $p$-divisibles.\piste{100-neron-extension-p-divisible}
\item Cote~130 : le symbole local de Néron se décompose en un nombre d'intersection sur le modèle de Néron et une correction caractérisée par quatre axiomes, d'où un accouplement à valeurs dans $\mathbf Z/m\mathbf Z$. Mathlib n'a pas les modèles de Néron ni les symboles locaux.\piste{130-neron-symbol-component-correction}
\item Cote~130 : les fonctions de Weil forment un faisceau de Zariski $QF_X$, extension de $\mathrm{Div}_X$ par $\Phi_X$ ; sur $X$ quasi-projectif, tout diviseur de Cartier est le diviseur d'une quasi-fonction globale. Mathlib n'a pas les fonctions de Weil ni les diviseurs de Cartier.\piste{130-quasi-function-sheaf}
\end{itemize}
